\documentclass[12pt]{article}
\usepackage[T1]{fontenc}
\usepackage[utf8]{inputenc}
\usepackage[brazil]{babel}
\usepackage[margin=1in]{geometry}
\usepackage{ragged2e}
\usepackage[normalem]{ulem}
\setlength{\parindent}{0pt}
\setlength{\parskip}{0.8\baselineskip}
\begin{document}
\justifying
\noindent Verifica-se que a questão objeto da controvérsia jurídica suscitada no recurso manejado pela parte recorrente esteve afetada no representativo da controvérsia - \textbf{Tema n. }\textbf{301} da TNU - (PU no processo n. PEDILEF 0501240-10.2020.4.05.8303/PE),  afetado em 17/03/2022 foi julgado (trânsito em julgado em 24/10/2022), por meio do qual foi elucidada a seguinte questão:

\noindent \textit{"Saber se, à luz da exigência de que o período de exercício de atividade rural seja imediatamente anterior ao requerimento de benefício ou implemento da idade, ainda que descontínuo, conforme arts. 39, i, 48, §2º e 143, todos da Lei 8.213/91, o exercício de atividade urbana por mais de 120 dias, corridos ou intercalados, no ano civil, na vigência da Lei 11.718/2008, implica, além da perda da qualidade de segurado especial, ruptura do perfil de trabalhador rural e interrupção da contagem do tempo de atividade rural (carência), impedindo o somatório dos períodos de atividade campesina anterior e posterior ao vínculo urbano que extrapolou o limite legal, exigindo nova contagem integral do intervalo exigido por lei para a aposentadoria por idade rural pura."}

\noindent Restou firmada a seguinte tese:

\noindent \textit{"Cômputo do Tempo de Trabalho Rural I. Para a aposentadoria por idade do trabalhador rural não será considerada a perda da qualidade de segurado nos intervalos entre as atividades rurícolas. Descaracterização da condição de segurado especial II. A condição de segurado especial é descaracterizada a partir do 1º dia do mês seguinte ao da extrapolação dos 120 dias de atividade remunerada no ano civil (Lei 8.213/91, art. 11, § 9º, III); III. Cessada a atividade remunerada referida no item II e comprovado o retorno ao trabalho de segurado especial, na forma do art. 55, parag. 3o, da Lei 8.213/91, o trabalhador volta a se inserir imediatamente no VII, do art. 11 da Lei 8.213/91, ainda que no mesmo ano civil." (Tema 301 TNU).}

\noindent Com efeito, depreende-se da leitura do voto condutor do acórdão que o mesmo se encontra em conformidade com o decidido no  tema pela TNU.

\noindent A proposito, consta do acordao hostilizado: (...) Aposentadoria por Tempo de Contribuição (Art. 55/6), Benefícios em Espécie, DIREITO PREVIDENCIÁRIO, Competência, Jurisdição e Competência, DIREITO PROCESSUAL CIVIL E DO TRABALHO

\noindent Nessas condições, o conteúdo do Acórdão recorrido é consentâneo com o entendimento da TNU, firmado em julgamento de recurso representativo de controvérsia, o que atrai a incidência da regra prevista pelo art. 14, III, \textit{b,} da Resolução CJF n. 586/2019:

\noindent \textit{\textbf{Art. 14.}}\textit{ Decorrido o prazo para contrarrazões, os autos serão conclusos ao magistrado responsável pelo exame preliminar de admissibilidade, que deverá, de forma sucessiva: (...) }\textit{\textbf{III }}\textit{- negar seguimento a pedido de uniformização de interpretação de lei federal interposto contra acórdão que esteja em conformidade com entendimento consolidado: }\textit{\textbf{b)}}\textit{ em }\uline{\textit{\textbf{recurso representativo de controvérsia pela Turma Nacional de Uniformização}}}\textit{ ou em pedido de uniformização de interpretação de lei dirigido ao Superior Tribunal de Justiça;(...).}

\noindent Posto isso, com arrimo no \uline{\textbf{art. 14, III, }}\uline{\textit{\textbf{b}}}\textit{,} da Resolução CJF n. 586/2019 c/c art. 1.030, I, do CPC, \textbf{NEGO SEGUIMENTO }\textbf{a}o\textbf{ PU}\textbf{.}

\noindent Intimem-se. Aguarde-se o decurso do prazo recursal. Após, certifique-se o trânsito em julgado e devolva-se ao juízo de origem.
\end{document}
